\begin{table}[t]
\centering
\caption{National County Preservation Summary}
\label{tab:county_preservation_summary}
\begin{tabular}{lr}
\toprule
\textbf{Metric} & \textbf{Value} \\
\midrule
Total States & 35 \\
Total Counties Analyzed & 2,636 \\
\midrule
Average Enacted Split Rate & 27.4\% \\
Average Algorithmic Split Rate & 28.4\% \\
Average Difference (Alg - Enacted) & +1.0\% \\
\midrule
Algorithmic Splits More Counties & 21 states \\
Algorithmic Splits Fewer Counties & 13 states \\
Similar Split Rates & 1 states \\
\midrule
\textbf{Category Classification} & \\
\quad Compactness-County Tradeoff & 4 states \\
\quad Algorithmic Preserves Better & 12 states \\
\quad Similar Preservation & 8 states \\
\midrule
Correlation (Splits $\leftrightarrow$ Compactness) & -0.68 \\
\bottomrule
\end{tabular}
\begin{tablenotes}
\small
\item \textit{Split Rate}: Percentage of counties split across multiple districts. \textit{Difference}: Positive means algorithmic splits more counties. \textit{Compactness-County Tradeoff}: States where algorithmic splits $>$3 more counties and gains $>$5\% compactness. \textit{Algorithmic Preserves Better}: States where algorithmic splits $<$3 fewer counties. \textit{Correlation}: Negative correlation ($-0.68$) indicates compactness gains often require splitting more counties.
\end{tablenotes}
\end{table}
